\clearpage
\begin{examplebox}
\textbf{\textcolor{CExample}{例 1（椭圆：已知中点求弦方程）}}

\textbf{题目：}已知椭圆 $\dfrac{x^2}{4}+\dfrac{y^2}{3}=1$，求以 $M(1,1)$ 为中点的弦所在直线方程。

\textbf{解：}先确认 $\dfrac{1^2}{4}+\dfrac{1^2}{3}=\dfrac7{12}<1$，中点在椭圆内部，非退化弦存在。
由 $a^2=4,b^2=3$，
\[
k=-\frac{b^2x_0}{a^2y_0}=-\frac34,
\quad y-1=-\frac34(x-1).
\]
\textbf{答案：}$\boxed{3x+4y-7=0}$。

\medskip
\textbf{\textcolor{CExample}{例 2（双曲线：验两实交点）}}

\textbf{题目：}已知双曲线 $\dfrac{x^2}{4}-\dfrac{y^2}{5}=1$，求以 $M(4,3)$ 为中点的弦所在直线方程。

\textbf{解：}先算 $Q=\dfrac{16}{4}-\dfrac95=\dfrac{11}{5}>1$，满足存在条件。
再用 $a^2=4,b^2=5$，
\[
k=\frac{5\cdot4}{4\cdot3}=\frac53,
\quad y-3=\frac53(x-4).
\]
\textbf{答案：}$\boxed{5x-3y-11=0}$。
代回双曲线得到 $55x^2-440x+664=0$，判别式
$\Delta=47520>0$ 且根的平均数为 $4$，确为所求实弦。

\tcbbreak
\textbf{\textcolor{CExample}{例 3（抛物线：斜率已知求中点轨迹）}}
\begin{center}\begin{tikzpicture}[font=\small,>=stealth]
\draw[gray!45,rounded corners=2pt] (0,0) rectangle (14.2,11.4);
\node at (7.1,10.75) {中点满足 $y_0=1$，但只落在一条开射线上};
\node at (7.1,9.9) {在 $y^2=4x$ 上画三条斜率均为 $2$ 的实弦。};
\begin{scope}[shift={(3.2,4)},scale=1]
\draw[gray,->] (-.5,0)--(5.3,0) node[right] {$x$};
\draw[gray,->] (0,-2.2)--(0,4.4) node[above] {$y$};
\draw[blue!75!black,thick,domain=-2.1:4.2,samples=100,variable=\v]
 plot ({\v*\v/4},\v);
\foreach \m/\col in {.75/orange!85!black,1.5/purple!80!black,2.5/teal!80!black}{
 \pgfmathsetmacro{\v}{sqrt(4*\m-1)}
 \draw[\col,thick] ({\m-\v/2},{1-\v})--({\m+\v/2},{1+\v});
 \fill[\col] ({\m-\v/2},{1-\v}) circle(1.6pt);
 \fill[\col] ({\m+\v/2},{1+\v}) circle(1.6pt);
 \fill[purple!80!black] (\m,1) circle(2.3pt);
}
\draw[purple!80!black,very thick,->] (.25,1)--(5.2,1);
\filldraw[fill=white,draw=purple!80!black,thick] (.25,1) circle(2.5pt)
 node[left=8pt,text=black] {$T$};
\end{scope}
\node[text=blue!75!black] at (10.9,8.2) {$y^2=4x$};
\node[text=purple!80!black] at (10.9,6.8) {中点在 $y=1$ 上移动};
\node at (10.9,5.6) {$T=(1/4,1)$ 为空心点};
\node at (10.9,4.5) {到 $T$ 时两端点重合};
\node at (10.9,3.2) {三条弦都满足 $k=2$};
\node at (7.1,1.15) {$\displaystyle \boxed{y_0=1,\quad x_0>\frac14}\qquad$切点边界不取，左侧也不可达。};
\end{tikzpicture}
\end{center}

\textbf{题目：}斜率为 $2$ 的直线与抛物线 $y^2=4x$ 交于不同实点 $A,B$，求中点 $M(x_0,y_0)$ 的轨迹。

\textbf{解：}在 $y^2=2px$ 中，$p=2$，故
\[
k=\frac{p}{y_0}\quad\Longrightarrow\quad
2=\frac2{y_0}\quad\Longrightarrow\quad y_0=1.
\]
\textbf{不能停在 $y=1$：}令两端点纵坐标为 $1+t$ 与 $1-t$（$t>0$），横坐标为
$\dfrac{(1+t)^2}{4},\dfrac{(1-t)^2}{4}$，因此
\[
x_0=\frac{1+t^2}{4}>\frac14.
\]
反之每个 $x_0>1/4$ 都可取 $t=\sqrt{4x_0-1}$ 实现。
\textbf{答案：}$\boxed{\ y=1,\ x>\dfrac14\ }$，不是整条直线。

\tcbbreak
\textbf{\textcolor{CExample}{例 4（高价值反例：求出斜率却没有弦）}}

\textbf{题目：}双曲线 $\dfrac{x^2}{4}-\dfrac{y^2}{5}=1$，是否存在以 $N(1,1)$ 为中点的实弦？

\textbf{解：}点差法先给出候选 $k=\dfrac{5}{4}$，即
$y-1=\dfrac54(x-1)$；但
\[
Q(N)=\frac14-\frac15=\frac1{20}\in(0,1).
\]
不满足 $Q<0$ 或 $Q>1$，\textbf{没有这样的非退化实弦}。
事实上代入直线后判别式为负（或直接用上述存在性判据）。
\textbf{启示：}“由中点算出斜率”\textbf{不是}“存在符合条件的弦”的充分证明。

\medskip
\textbf{\textcolor{CExample}{例 5（分母为零：中心弦不能套公式）}}

\textbf{题目：}椭圆 $\dfrac{x^2}{4}+y^2=1$ 中，若弦的中点是 $O(0,0)$，能否断定弦为竖直弦？

\textbf{解：}不能。$y=0$ 交椭圆于 $(-2,0),(2,0)$，是水平中心弦；
$x=0$ 交椭圆于 $(0,-1),(0,1)$，是竖直中心弦。
\textbf{结论：}中心对应无穷多条直径方向；此时 $0/0$ 型公式不能使用。

\medskip
\textbf{\textcolor{CExample}{变式（考场迁移）}}

对圆 $x^2+y^2=25$，已知中点 $M(3,0)$，求弦所在直线。
此时不能代 $k=-x_0/y_0$；由圆的对称性或无分母点差式
$3(x-3)=0$ 得 $\boxed{x=3}$，交圆于 $(3,\pm4)$。
\end{examplebox}
